\documentclass[10pt, letterpaper]{article}

% Packages:
\usepackage[
    ignoreheadfoot, % set margins without considering header and footer
    top=2 cm, % seperation between body and page edge from the top
    bottom=2 cm, % seperation between body and page edge from the bottom
    left=2 cm, % seperation between body and page edge from the left
    right=2 cm, % seperation between body and page edge from the right
    footskip=1.0 cm, % seperation between body and footer
    % showframe % for debugging 
]{geometry} % for adjusting page geometry
\usepackage{titlesec} % for customizing section titles
\usepackage{tabularx} % for making tables with fixed width columns
\usepackage{array} % tabularx requires this
\usepackage[dvipsnames]{xcolor} % for coloring text
\definecolor{primaryColor}{RGB}{0, 0, 0} % define primary color
\usepackage{enumitem} % for customizing lists
\usepackage{fontawesome5} % for using icons
\usepackage{amsmath} % for math
\usepackage[
    pdftitle={Flavio Romano's CV},
    pdfauthor={Flavio Romano},
    colorlinks=true,
    urlcolor=primaryColor
]{hyperref} % for links, metadata and bookmarks
\usepackage[pscoord]{eso-pic} % for floating text on the page
\usepackage{calc} % for calculating lengths
\usepackage{bookmark} % for bookmarks
\usepackage{lastpage} % for getting the total number of pages
\usepackage{changepage} % for one column entries (adjustwidth environment)
\usepackage{paracol} % for two and three column entries
\usepackage{ifthen} % for conditional statements
\usepackage{needspace} % for avoiding page brake right after the section title
\usepackage{iftex} % check if engine is pdflatex, xetex or luatex

% Ensure that generate pdf is machine readable/ATS parsable:
\ifPDFTeX
    \input{glyphtounicode}
    \pdfgentounicode=1
    \usepackage[T1]{fontenc}
    \usepackage[utf8]{inputenc}
    \usepackage{lmodern}
\fi

\usepackage{charter}

% Some settings:
\raggedright
\AtBeginEnvironment{adjustwidth}{\partopsep0pt} % remove space before adjustwidth environment
\pagestyle{empty} % no header or footer
\setcounter{secnumdepth}{0} % no section numbering
\setlength{\parindent}{0pt} % no indentation
\setlength{\topskip}{0pt} % no top skip
\setlength{\columnsep}{0.15cm} % set column seperation
\pagenumbering{gobble} % no page numbering

\titleformat{\section}{\needspace{4\baselineskip}\bfseries\large}{}{0pt}{}[\vspace{1pt}\titlerule]

\titlespacing{\section}{
    % left space:
    -1pt
}{
    % top space:
    0.3 cm
}{
    % bottom space:
    0.2 cm
} % section title spacing

\renewcommand\labelitemi{$\vcenter{\hbox{\small$\bullet$}}$} % custom bullet points
\newenvironment{highlights}{
    \begin{itemize}[
        topsep=0.10 cm,
        parsep=0.10 cm,
        partopsep=0pt,
        itemsep=0pt,
        leftmargin=0 cm + 10pt
    ]
}{
    \end{itemize}
} % new environment for highlights


\newenvironment{highlightsforbulletentries}{
    \begin{itemize}[
        topsep=0.10 cm,
        parsep=0.10 cm,
        partopsep=0pt,
        itemsep=0pt,
        leftmargin=10pt
    ]
}{
    \end{itemize}
} % new environment for highlights for bullet entries

\newenvironment{onecolentry}{
    \begin{adjustwidth}{
        0 cm + 0.00001 cm
    }{
        0 cm + 0.00001 cm
    }
}{
    \end{adjustwidth}
} % new environment for one column entries

\newenvironment{twocolentry}[2][]{
    \onecolentry
    \def\secondColumn{#2}
    \setcolumnwidth{\fill, 4.5 cm}
    \begin{paracol}{2}
}{
    \switchcolumn \raggedleft \secondColumn
    \end{paracol}
    \endonecolentry
} % new environment for two column entries

\newenvironment{threecolentry}[3][]{
    \onecolentry
    \def\thirdColumn{#3}
    \setcolumnwidth{, \fill, 4.5 cm}
    \begin{paracol}{3}
    {\raggedright #2} \switchcolumn
}{
    \switchcolumn \raggedleft \thirdColumn
    \end{paracol}
    \endonecolentry
} % new environment for three column entries

\newenvironment{header}{
    \setlength{\topsep}{0pt}\par\kern\topsep\centering\linespread{1.5}
}{
    \par\kern\topsep
} % new environment for the header

\newcommand{\placelastupdatedtext}{% \placetextbox{<horizontal pos>}{<vertical pos>}{<stuff>}
  \AddToShipoutPictureFG*{% Add <stuff> to current page foreground
    \put(
        \LenToUnit{\paperwidth-2 cm-0 cm+0.05cm},
        \LenToUnit{\paperheight-1.0 cm}
    ){\vtop{{\null}\makebox[0pt][c]{
        \small\color{gray}\textit{Last updated in September 2024}\hspace{\widthof{Last updated in September 2024}}
    }}}%
  }%
}%

% save the original href command in a new command:
\let\hrefWithoutArrow\href

% new command for external links:


\begin{document}
\newcommand{\AND}{\unskip
	\cleaders\copy\ANDbox\hskip\wd\ANDbox
	\ignorespaces
}
\newsavebox\ANDbox
\sbox\ANDbox{$|$}

\begin{header}
	\fontsize{25 pt}{25 pt}\selectfont Flavio Romano
	
	\vspace{5 pt}
	
	\normalsize
	\mbox{Italy}%
	\kern 5.0 pt%
	\AND%
	\kern 5.0 pt%
	\mbox{\hrefWithoutArrow{mailto:flavromano@proton.me}{flavromano@proton.me}}%
	\kern 5.0 pt%
	\AND%
	\kern 5.0 pt%
	\mbox{\hrefWithoutArrow{https://www.linkedin.com/in/flavromano/}{linkedin.com/in/flavromano}}%
	\kern 5.0 pt%
	\AND%
	\kern 5.0 pt%
	\mbox{\hrefWithoutArrow{https://github.com/FlavRomano}{github.com/FlavRomano}}%
\end{header}

\vspace{5 pt - 0.3 cm}

\vspace{.4 cm}

Software Engineer and Technical Lead specializing in backend systems, IAM solutions, and scalable cloud-native architectures. Experienced in leading cross-functional teams, designing distributed systems, and delivering enterprise-grade software in telecommunications and sustainability domains.

Currently pursuing an MSc in Artificial Intelligence while working part-time at NTT DATA, strengthening a software engineering background with advanced AI and machine learning expertise. Passionate about applying AI to solve complex industrial problems at scale.
\section{Experience}

\begin{twocolentry}{
		\textbf{March 2024 – Present}
	}
	\textbf{Software Engineer}, NTT DATA -- Pisa, IT
\end{twocolentry}

\vspace{0.10 cm}
\begin{onecolentry}
	
	\textbf{Identity \& Access Management Platform (Telco Client)}
	\hfill Aug 2024 -- May 2025
	
	\begin{highlights}
		\item Contributed to the design and implementation of OAuth 2.0, OpenID Connect, and SAML-based authentication solutions for a leading Italian telecommunications provider.
		
		\item Enhanced a customized Keycloak deployment, enabling adoption of the Phantom Token architecture to improve security, scalability, and operational flexibility.
		
		\item Developed Spring Boot microservices supporting authorization workflows, including token introspection, redirect management, and third-party JWT/JWE processing.
		
		\item Collaborated within a 5-person engineering team delivering authentication and identity management capabilities across multiple enterprise portals.
		
	\end{highlights}
	
	\vspace{0.15 cm}
	
	\textbf{Logistics Data Processing Platform}
	\hfill May 2025 -- Sep 2025
	
	\begin{highlights}
		\item Designed and implemented Spring Batch jobs for large-scale shipment processing, handling hundreds of thousands of records through complex transformation and enrichment pipelines.
		
		\item Developed resilient batch-processing workflows with a focus on reliability, maintainability, and performance optimization.
		
	\end{highlights}
	
	\vspace{0.15 cm}
	
	\textbf{AI-powered Due Diligence Platform for Sustainability Compliance}
	\hfill Sep 2025 -- Jan 2026
	
	\begin{highlights}
		\item Served as Technical Lead for the development of an AI-powered platform designed to simplify the generation of due diligence documentation required by European import/export sustainability regulations.
		
		\item Coordinated a team of 4 engineers, managing task allocation, technical mentoring, and knowledge-transfer sessions.
		
		\item Designed and implemented backend services using Spring Boot, Redis caching, and optimized offset- and cursor-based pagination strategies to support large datasets efficiently.
		
		\item Worked closely with stakeholders to translate regulatory and business requirements into scalable software solutions.
		
	\end{highlights}
	
	\vspace{0.15 cm}
	
	\textbf{Corporate Sustainability \& CO$_2$ Emissions Assessment Platform}
	\hfill Jan 2026 -- Present
	
	\begin{highlights}
		\item Acting as Technical Lead for the development of a platform that estimates CO$_2$ emissions generated by software systems, hardware infrastructure, and AI workloads.
		
		\item Leading a team of 5 engineers, coordinating daily stand-ups, milestone planning, task assignment, code reviews, and technical decision-making.
		
		\item Serving as the primary technical point of contact for the client, gathering requirements and aligning technical deliverables with business objectives.
		
		\item Designing scalable backend services and data-processing workflows to support sustainability reporting and environmental impact assessment.
		
	\end{highlights}
	
\end{onecolentry}

\vspace{0.2 cm}

\begin{twocolentry}{
		\textbf{Dec 2023 – March 2024}
	}
	\textbf{Software Engineer Stage}, NTT DATA -- Pisa, IT
\end{twocolentry}

\vspace{0.10 cm}
\begin{onecolentry}
	\begin{highlights}
		\item Developed a Java microservice for the GreenIT ecosystem to collect hardware metrics and estimate Software Carbon Intensity (SCI) and CO$_2$ emissions on Google Cloud Platform infrastructures.
		
	\end{highlights}
\end{onecolentry}


\section{Technical Skills}
\begin{onecolentry}
	\textbf{Programming languages:} Bash, C, Golang, Java, JS/TS, Python, SQL
\end{onecolentry}
        
\vspace{0.2 cm}

\begin{onecolentry}
	\textbf{Framework:} Gin, Pandas, Selenium, Spring, SvelteKit, Pytorch, Keras
\end{onecolentry}

\vspace{0.2 cm}

\begin{onecolentry}
	\textbf{Cloud provider:} AWS, GCP
\end{onecolentry}

\section{Projects}      

\begin{twocolentry}{
		Apr 2026 – May 2026
	}
	\textbf{Encoding Classification using LSTM} (\href{https://github.com/FlavRomano/encoding-classifier-lstm}{github})
\end{twocolentry}

\begin{highlights}
	\item Developed a PyTorch-based sequence classification model to automatically identify the encoding algorithm applied to a text string, distinguishing between Caesar Cipher, ROT13, and MD5 transformations.
	
	\item Built and trained an LSTM neural network with approximately 230k trainable parameters using a dataset of more than 72,000 samples derived from English-language anagrams.
	
	\item Designed an NLP preprocessing pipeline based on token embeddings and sequence padding, achieving 97.99% test accuracy on unseen data.
	
	\item Implemented automated training and inference workflows, enabling real-time prediction of encoding types from arbitrary input strings.
\end{highlights}


\begin{twocolentry}{
		Oct 2025 – Jan 2026
	}
	\textbf{Multi-output Regression for Noisy Datasets}
\end{twocolentry}

\vspace{0.10 cm}
\begin{onecolentry}
	\begin{highlights}
		\item Developed a deep learning pipeline for a multi-output regression task using a highly noisy dataset. Achieved significant error reduction on the test set through rigorous data preprocessing and feature engineering.
		\item Conducted systematic model selection and hyperparameter tuning using Optuna, comparing various Neural Network architectures implemented with Keras and Scikit-learn.
		\item Performed a robust final model assessment to ensure generalization, earning high marks for the methodological approach in the Machine Learning graduate course.
	\end{highlights}
\end{onecolentry}
\vspace{0.2 cm}
\begin{twocolentry}{
		Aug 2023 - Dec 2023
	}
	\textbf{Coinquilining} (\href{https://github.com/FlavRomano/coinquilining}{github})\end{twocolentry}
	
	\vspace{0.10 cm}
	\begin{onecolentry}
		\begin{highlights}
			\item Developed from scratch as final project for the "Web App Development" course, earning a top grade (30/30). Coinquilining is a PWA for managing roommate life, featuring shared fridge and pantry management, expense tracking, and recipe suggestions.
			\item Implemented expiration tracking with notifications, ingredient-based recipe searches using Puppeteer to scrape a dynamic site, and a calendar for fair task distribution (time shift problem).
			\item Tech Stack: Supabase, TypeScript, SvelteKit, Tailwind CSS, Puppeteer
		\end{highlights}
	\end{onecolentry}
	
	\section{Education}
	
	\begin{twocolentry}{
			Sept 2025 – Present
		}
		\textbf{University of Pisa}, Master of Science (MSc) in Artificial Intelligence \end{twocolentry}
		
		\vspace{0.10 cm}
		\begin{twocolentry}{
				Sept 2020 – May 2024
			}
			\textbf{University of Pisa}, Bachelor of Computer Science\end{twocolentry}
			
			\vspace{0.10 cm}
			\begin{onecolentry}
				\begin{highlights}
					\item \textbf{Thesis:} Development of a configurable connector to monitoring SCI and CO2 emissions of GreenIT® Google Cloud Platform components.
				\end{highlights}
			\end{onecolentry}
			
\end{document}